\documentclass[a4paper]{article}

\usepackage{graphicx}
\usepackage{amsmath,amssymb}
\usepackage{minted}
\usepackage{multirow}
\usepackage{float}
\usepackage{algorithm}
\usepackage{algpseudocode}
\usepackage{todonotes}
\usepackage{forest}
\usepackage{xstring}
\usepackage{xcolor}
\usepackage[most]{tcolorbox}
\usepackage{xparse}
\usepackage{natbib}
\usepackage{adjustbox}
\usepackage[active,tightpage]{preview}

\PreviewBorder=4pt

\usetikzlibrary{graphs,graphdrawing,arrows.meta,positioning,calc,fit,shapes.geometric}
\usegdlibrary{layered}

% for external graphviz
\input{/nd_language_ai_project/src/nd_language_ai/formal/computer/programming/tex/latex/code_clips/external_graphviz.tex}

% for tags
\input{/nd_language_ai_project/src/nd_language_ai/formal/computer/programming/tex/latex/part/control_sequence/kind/semtag/semtag.tex}

% for links
\input{/nd_language_ai_project/src/nd_language_ai/formal/computer/programming/tex/latex/code_clips/seehere.tex}

% for nested lists and sections
\input{/nd_language_ai_project/src/nd_language_ai/formal/computer/programming/tex/latex/code_clips/deep_structure.tex}

\title{}
\date{}
\author{Mohammad Rahmani}

\begin{document}
   \maketitle

   \seeherefooter{https://en.wikipedia.org/wiki/Free_energy_principle#Active_inference}
   \url{https://en.wikipedia.org/wiki/Free_energy_principle}
   \includepdf[pages=-, fitpaper=true]{\assetspath review/biology/active_inference.pdf}
   \seeherefooter{}
   Active inference can explain cognitive proccesses such as perception, action selection and

   \bibliographystyle{plainnat}
   \bibliography{/home/donkarlo/Dropbox/repo/nd_shared_research/bibliography/bibliography.bib}
\end{document}
